% Results prose. Numbers come from numbers.tex macros where one exists; the few
% literal values are from results/tables/sensitivity.csv and failure_*.csv.

\newcommand{\RESULTSMAIN}{%
Over the full out-of-sample span the MTL strategy earns an annualized Sharpe ratio of \MtlAllSharpe{} (CAGR \MtlAllCagr, maximum drawdown \MtlAllMdd, weekly win rate \MtlAllWin, profit factor \MtlAllPf) with an information ratio of \MtlAllIr{} against the Nifty~50, which returned \NiftyAllCagr{} p.a.\ with a Sharpe of \NiftyAllSharpe{} over the same weeks. Every baseline that trades the same universe through the same engine is close to zero: momentum \MomAllSharpe, inverse volatility \InvVolAllSharpe, random 5-stock portfolios \RandAllSharpe. The graph-feature ablation retains part of the edge (\NoGraphAllSharpe) but with a deeper drawdown (\NoGraphAllMdd). Two caveats are visible in the same table. First, the 95\% block-bootstrap interval for the MTL Sharpe, \MtlAllSharpeCi, includes zero: 167 weeks of a five-stock portfolio cannot establish significance on their own. Second, the mean weekly rank-IC over the whole span is only \MtlAllIc, which is far smaller than the portfolio statistics suggest---the two disagree because, as the period breakdown shows next, the signal is concentrated in time.}

\newcommand{\RESULTSPERIODS}{%
Table~\ref{tab:periods} splits the span into the PRD test period (2023-H2 to 2024, \MtlTestWeeks{} weeks) and the untouched holdout (2025 to September~2026, \MtlHoldWeeks{} weeks). The two halves could hardly differ more. In the test period the strategy earns a Sharpe of \MtlTestSharpe{} (CAGR \MtlTestCagr, drawdown \MtlTestMdd, win rate \MtlTestWin, IC \MtlTestIc) while the Nifty~50 returned \NiftyTestCagr{} (Sharpe \NiftyTestSharpe); the quarterly rank-IC of the return head was 0.10, 0.10, 0.15 and 0.10 in the four quarters from 2023-Q3 to 2024-Q2, an unusually strong cross-sectional signal. In the holdout the strategy \emph{loses} money: Sharpe \MtlHoldSharpe, CAGR \MtlHoldCagr, IC \MtlHoldIc, against a flat market (Nifty CAGR \NiftyHoldCagr, Sharpe \NiftyHoldSharpe) in which momentum (\MomHoldSharpe) and random portfolios (\RandHoldSharpe) also fail. The calendar-year view in Table~\ref{tab:yearly} makes the transition explicit: the model beat the index by 31 and 51 percentage points in 2023-H2 and 2024, then trailed it by 15 points in 2025 and by 5 points in 2026. The rolling IC (Fig.~\ref{fig:rollic}) crosses zero in the third quarter of 2024 and never recovers. We return to why in Sec.~\ref{sec:limits}; the decile plot of Fig.~\ref{fig:decile} shows the pooled outcome: the top decile earns the highest mean weekly return, but the pattern across deciles is not monotone.}

\newcommand{\RESULTSABLATION}{%
Table~\ref{tab:ablation} compares the full model with its ablations on the PRD test period, where all variants were run (the return-only and cost-term variants were not extended into the holdout for compute reasons; the no-graph variant was). Removing the graph features lowers the test Sharpe from \MtlTestSharpe{} to \NoGraphTestSharpe{} and the IC from \MtlTestIc{} to \NoGraphTestIc; over the full span the gap is \MtlAllSharpe{} versus \NoGraphAllSharpe. Removing the volatility task is more damaging: a return-only model with the same backbone reaches only \RetOnlyTestSharpe{} (IC \RetOnlyTestIc) and trades far more (turnover 0.92 versus 0.36 per week), which is the clearest evidence in the study that the auxiliary volatility task regularizes the shared representation. Adding the transaction-cost term ($\lambda=0.3$) to the tuned configuration did \emph{not} help (\CostTestSharpe, IC \CostTestIc, turnover 0.69): the tuner had already rejected $\lambda>0$ on validation IC, and the softmax-implied portfolio used by the penalty is a poor proxy for the hard top-decile portfolio that is actually traded. The logistic-regression baseline on the same last-week and graph features is a strong competitor on the test period (\LogitTestSharpe, IC \LogitTestIc) at higher turnover (0.86); we did not run it on the holdout. The static single-fit models tell the same story in IC terms (val / test / holdout): full model \staticIcFinalVal{} / \staticIcFinalTest{} / \staticIcFinalHoldout; no-graph \staticIcFinalnographVal{} / \staticIcFinalnographTest{} / \staticIcFinalnographHoldout; return-only \staticIcFinalretonlyVal{} / \staticIcFinalretonlyTest{} / \staticIcFinalretonlyHoldout; and selecting the stopping epoch on validation Sharpe instead of IC gives \staticIcFinalsharpestopVal{} / \staticIcFinalsharpestopTest{} / \staticIcFinalsharpestopHoldout.}

\newcommand{\RESULTSSENS}{%
Table~\ref{tab:sens} and Fig.~\ref{fig:cost} vary the backtest assumptions on the full span. The strategy survives realistic costs: the Sharpe falls from 0.90 with free trading to 0.77 at the 0.1\%/side base case and 0.57 at 0.25\%/side, a slope consistent with its 0.47 average weekly one-way turnover. The hold-zone rule cuts turnover to 0.18 at a modest cost in Sharpe (0.67) and improves the drawdown (\(-23.4\%\)). A dollar-neutral long--short book is worse (0.40): the SELL decile carries little information beyond what the BUY decile already exploits, and shorting doubles the cost burden. Widening the BUY set to the top 20\% dilutes the edge (0.55). The 5\% position cap prescribed in the PRD leaves 75\% of capital in cash and therefore only rescales the CAGR (10.3\% versus 20.9\%); the Sharpe is unchanged because cash earns the risk-free rate. The last block of the table is a post-hoc diagnostic of the decision layer that we did not use for any selection: ranking on the predicted return alone would have delivered a Sharpe of 0.82 over the full span and \emph{flat} performance in the holdout (\(-0.11\), IC \(+0.014\)) instead of the \(-1.12\) obtained with $\hat r/\hat\sigma$; a shrunk ratio $\hat r/(\hat\sigma+0.03)$ sits in between (0.95 full span, \(-0.21\) holdout). Ranking on predicted volatility alone is strongly negative (\(-0.80\)). In other words, the Sharpe-style score is the right idea when the return forecast is informative, but when that forecast decays towards zero, dividing by volatility turns noise into a systematic tilt.}

\newcommand{\RESULTSATT}{%
The attention analysis produced a negative result that we consider informative. In the tuned model both the self-attention and the pooling weights are exactly uniform (entropy \(4.094\) versus \(\log 60 = 4.094\); the share on the last four weeks is \attLastFour, i.e.\ $4/60$). Inspection of the weights shows why: the selected weight decay (\bestweightdecay, roughly 57 times the PRD default of $10^{-5}$) shrinks the attention projection matrices to a norm of $10^{-3}$, so the sequence encoder degenerates into a mean over the 60 tokens followed by the feed-forward layers, and the model effectively becomes ``pooled own-history $+$ graph features''. The PRD-default configuration (weight decay $10^{-5}$) retains a non-trivial attention pattern (projection norm 10.0), shown in Fig.~\ref{fig:att}: the most recent 3--4 weeks and the oldest positions receive above-uniform attention in the last layer, and the pooling weights tilt towards the recent weeks in high-volatility regimes and away from them in calm regimes, although the magnitudes are small (maximum 0.021 against a uniform 0.017). The lag profile of SHAP mass (Fig.~\ref{fig:shaplag}) is likewise nearly flat (the last four weeks carry 5.7\% of sequence attribution, the oldest twelve 21\%). The honest summary is that, at this data size ($\approx$9{,}500 training windows), the optimizer found no value in modelling \emph{which} past weeks matter and the validation-selected model is closer to a well-regularized MLP than to a transformer.}

\newcommand{\RESULTSSHAP}{%
SHAP attributions of the decision score on 600 test-period samples (Fig.~\ref{fig:shap}) are dominated by graph features: fourteen of the fifteen most important inputs are graph features, led by \topFeatureOne, \topFeatureTwo{} and \topFeatureThree, and the graph groups together account for \shapGraphShare{} of attribution mass (67\% if sequence attributions are aggregated by absolute value over lags rather than net). The correlation-trend group (betas, market and sector correlations) is the largest single group, followed by the market-regime descriptors and peer momentum. Importance is spread out: the five largest features explain \shapTopFive{} of the total and the ten largest \shapTopTen, so the PRD's target of ``top five features explain 60\%'' is not met. The direction of the largest effects is consistent with the failure analysis below: high market beta and high market correlation raise the score, i.e.\ the model learned to buy high-beta, market-sensitive names, which is exactly what worked in the 2023--2024 rally and stopped working afterwards.}

\newcommand{\RESULTSFAIL}{%
Table~\ref{tab:regimes} and Fig.~\ref{fig:regimes} condition the out-of-sample weeks on regime variables. The signal is clearly better when pairwise correlations are low (IC 0.039, $t=1.8$) than when they are high (\(-0.024\)), and in rising markets (12-week trend above median: IC 0.028, portfolio Sharpe 1.50) than in falling ones (\(-0.012\), 0.12); market volatility and dispersion make little difference. In the 19 weeks in which the index fell by more than 2\% the cross-sectional IC is actually high (0.091) but the long-only portfolio loses heavily (mean \(-2.5\%\) per week), i.e.\ the model ranks stocks sensibly inside a crash but cannot avoid the crash. By stock characteristic, BUY signals on high-beta names had the best hit rate (57\%, mean next-week return \(+1.1\%\)) against 46\% for low-beta names; IT stocks were the worst sector (hit rate 40\%, mean \(-1.3\%\)) and financial services the best (58\%). A final failure mode is procedural: with a 26-week validation window the IC-based early-stopping criterion sometimes selected the first epoch (in 4 of 39 seed-fold fits), so several folds are trained models in name only; seed averaging mitigates but does not remove this.}

\newcommand{\RESULTSDISC}{%
\textbf{What worked.} The volatility head is robust: its rank-IC against realized next-week volatility stays between 0.18 and 0.38 in every quarter, including the holdout (mean 0.26), and it beats the naive ``last week's volatility'' forecast in every split. Its presence improves the return head (return-only ablation) and halves turnover. The graph features carry most of the attribution mass and their removal costs performance in every period. And the system as specified---weekly decile portfolio, 0.1\% costs, walk-forward retraining---delivered a large, cost-robust edge for eighteen months.

\textbf{What did not.} The return signal is regime-bound. Its quarterly IC exceeded 0.10 throughout the 2023--2024 rally, during which the model's picks carried twice the 12-week momentum of the universe, and it fell to zero or below from 2024-Q3 onward while the model's picks lost their momentum tilt; the 2025--2026 holdout, a flat-to-down market, is where a strategy built on beta, correlation and momentum features should be expected to fail, and it did. The decision layer amplified the failure: dividing a near-zero return forecast by predicted volatility produced a systematic tilt with a significantly negative IC, whereas ranking on the return alone would have been flat. We report the holdout precisely because it is where the PRD targets (Sharpe $>1$, win rate $>52\%$) are missed; on the PRD test period alone they are exceeded by a wide margin (\MtlTestSharpe, \MtlTestWin).

\textbf{Methodological lessons.} (i)~Selecting on a 26-week validation Sharpe is close to selecting on noise (Fig.~\ref{fig:curves}); rank-IC is better but still weak enough to pick epoch~1 in some folds, and a longer validation window or a smoothed criterion should be used. (ii)~Hyper-parameter search on a short window can silently disable a component: the tuned weight decay removed the attention path, so the architecture that is reported is not the architecture that was trained. (iii)~One diagnostic during development (an IC computation for an untuned model) did look at the 2023-H2--2024 period; the holdout was never inspected, and it is the only period we regard as a clean test. (iv)~The differentiable cost term as formulated did not transfer to the hard top-decile portfolio; a penalty on the change in \emph{ranks}, or direct optimization of the traded portfolio's net return, is the natural next step. (v)~A model of this kind should be paired with a regime gate---the low-correlation / rising-trend conditions under which the signal has positive IC are observable at decision time.}

\newcommand{\RESULTSCONC}{%
On the PRD test period the system exceeded every target (Sharpe \MtlTestSharpe, win rate \MtlTestWin, information ratio \MtlTestIr); on the untouched 2025--2026 holdout it lost money (Sharpe \MtlHoldSharpe) as its beta- and momentum-driven return signal decayed and the Sharpe-style decision rule amplified the resulting noise; over the full span it remains ahead of every baseline (Sharpe \MtlAllSharpe{} against \NiftyAllSharpe{} for the index) but not significantly so.}
